% !TeX root = ../../Book2.tex
\chapter{Bochner–Riesz 平均}
\label{b2:ch:24}
\BookChapterNumber{b2:ch:24}
\begin{chapterabstract}
Bochner–Riesz 平均以边界消失阶平滑球形频率截断。本章计算完整核常数与振荡渐近，证明可积核阈值、尾部必要条件及一个解析插值充分范围，并说明限制估计如何控制薄壳。Fefferman 球乘子否定定理作为准确标注的深层输入陈述。
\end{chapterabstract}

\begin{learninggoals}\begin{itemize}
\item\label{b2:item:24-goal-1}能核对平均的缩放、Bessel 核与原点常数。
\item\label{b2:item:24-goal-2}能从核渐近的非零下界推导严格必要条件。
\item\label{b2:item:24-goal-3}能区分可积核范围、已证插值范围与经典猜想范围。
\item\label{b2:item:24-goal-4}能将球面限制用于薄壳混合指数估计，避免把它误写成同指数乘子界。
\end{itemize}\end{learninggoals}

矩形硬截断可用各坐标 Hilbert 投影组合，球面硬截断则同时涉及连续变化的法向与曲率。加入正参数平滑后，符号边界消失的程度决定核的衰减。下面以核计算和已证范围为主线，把猜想与深层否定结果放在清楚的证明层次中。

% !TeX root = ../../Book2.tex
\section{Bochner–Riesz 乘子}
\label{b2:sec:2401}

球形频率截断与矩形截断具有不同的几何。Bochner–Riesz 平均在球边界加入消失因子，使硬截断变成带光滑度参数的族。问题是需要多少边界平滑，才能在给定 \(L^p\) 上得到一致范数界与逼近。

\subsection{定义}
\label{b2:subsec:240101}

对 \(\delta\ge0\)、\(R>0\)，定义
\[
 B_R^\delta f=\mathcal F^{-1}\!\left[
       (1-|\xi|^2/R^2)_+^\delta\widehat f(\xi)\right].
\]
\symidx{\(B_R^\delta f\)：Bochner--Riesz 平均}
在 \(\delta=0\) 时取球特征函数，边界值无关。先在 \(L^2\) 或 Schwartz 函数上定义，再问其 \(L^p\) 延拓。记单位半径核为 \(K_\delta\)。正参数 \(\delta\) 保证符号在边界连续，却不保证任意阶光滑；若 \(\delta\) 非整数，跨边界导数的阶数受到限制。

\ProseParagraphBreak
Bochner–Riesz 平均不是简单的曲面平均，它包含球内部的频率；但边界附近的薄壳是最难控制的部分。远离球边界的符号可用一个固定光滑截断处理，其核属于 Schwartz 空间，所有 \(L^p\) 上都无困难。

\subsection{缩放}
\label{b2:subsec:240102}

缩放给出
\[
 K_{\delta,R}(x)=R^dK_\delta(Rx),\qquad
 B_R^\delta=D_RB_1^\delta D_{1/R}.
\]
因此同指数 \(L^p\) 范数独立于 \(R\)，前提是单位半径算子确实在该空间有界。变换归一化中的 \(2\pi\) 保留在核公式中。

\ProseParagraphBreak
若已知 \(B_1^\delta\in M_p\)、\(1<p<\infty\)，则
\(B_R^\delta f\to f\) 在 \(L^p\) 中成立。证明先取频率紧支撑的 Schwartz 函数：符号在这个紧集上随着 \(R\to\infty\) 连同所需导数趋于 \(1\)，所以逆变换误差在 Schwartz 范数中趋零。这样的函数在 \(L^p\) 中稠密，统一算子界再把结论推广到所有输入。
这一论证不涉及逐点极大算子；范数收敛与几乎处处全参数收敛是不同问题。

% 扩写定位：b2:subsec:240102
平均阶数增加时，有界性具有一个精确的单调传递原理。
\begin{proposition}\label{b2:prop:24-order-monotonicity}
设 \(0\leq\delta_1<\delta_2\)、\(\eta=\delta_2-\delta_1\)。则
\[
 B_R^{\delta_2}f=
 \frac1{\mathrm B(\delta_1+1,\eta)}
 \int_0^1 B_{R\sqrt t}^{\delta_1}f\,
 t^{\delta_1}(1-t)^{\eta-1}\,\mathrm dt.
\]
若 \(\sup_R\|B_R^{\delta_1}\|_{p\to p}\leq C\)、\(1\leq p<\infty\)，则同一界对 \(\delta_2\) 成立。极大算子的点态上界也满足
\(\sup_R|B_R^{\delta_2}f|\leq\sup_R|B_R^{\delta_1}f|\)，只要右侧由共同的测试函数定义或已延拓的极大算子解释。
\end{proposition}
\begin{proof}
先对 Schwartz 输入在 Fourier 侧验证。令 \(a=|\xi|^2/R^2\)。当 \(a<1\)，右侧符号为
\[
 \frac1{\mathrm B(\delta_1+1,\eta)}
 \int_a^1(t-a)^{\delta_1}(1-t)^{\eta-1}\,\mathrm dt
 =(1-a)^{\delta_1+\eta},
\]
换元 \(t=a+(1-a)u\) 就得到 beta 积分。若 \(a\geq1\)，两边都为零。归一化权的总积分为一，因而 Minkowski 给出算子界，逐点绝对值给出极大界。有限 \(p\) 的连续延拓可由测试函数稠密性得到；在有界输入上需要明确采用核、弱对偶或其他已经指定的延拓，而不是依赖不存在的 Schwartz 范数稠密性。
\end{proof}
因此若某个阶数已给出一致 \(L^p\) 有界性，所有更高阶也有界。但这个结论只在阶数方向向上推，不能从可积核的大阶数反推临界附近的界。其权在 \(\eta\downarrow0\) 时集中于 \(t=1\)，没有提供将平滑程度降低的逆算子。

\ProseParagraphBreak

这个公式同时解释了“平均”的一层含义：高阶 Bochner--Riesz 是较低阶频率截断按半径作正权平均。正的是半径平均的权，并不意味着最后得到的实空间卷积核处处非负。频率截断自身已经有振荡核，正权平均也可能保留符号变化。

\ProseParagraphBreak

阶数与半径若同时变化，还会出现不同的极限模型。令 \(\delta_n\to\infty\)、\(R_n\to\infty\)，并假定 \(\delta_n/R_n^2\to a\in[0,\infty)\)。对每个固定频率，
\[
 (1-|\xi|^2/R_n^2)_+^{\delta_n}\longrightarrow e^{-a|\xi|^2}.
\]
这是对数展开的结果：\(\log(1-v)=-v+O(v^2)\)，而在比例有界时 \(\delta_n/R_n^4\to0\)。所有符号都介于零和一之间，故 Plancherel 与支配收敛给出在 \(L^2\) 上的强算子极限。\(a=0\) 是恒等；\(a>0\) 是时间 \(a/(4\pi^2)\) 的热算子。

\ProseParagraphBreak

若比例 \(\delta_n/R_n^2\to\infty\)，则每个非零频率上的符号由
\((1-v)^{\delta_n}\leq e^{-\delta_n v}\) 控制而趋零，同样得到 \(L^2\) 强极限零算子。零频单点不影响平方密度。因此“半径趋于无穷便逼近恒等”需要阶数固定或至少记录其增长；两个参数改变时不能忽略彼此的尺度关系。

\ProseParagraphBreak

特别取 \(R=\sqrt\delta\)，就由球内多项式平均得到 Gaussian 热模型。这是在 Fourier 侧的平方极限，不是每个有限阶实空间核已经为正的断言。有限阶 Bessel 核仍有正负振荡；极限的正性与近似对象逐个正性是不同问题。类似地，平方收敛本身没有证明全体 \(L^p\) 的统一算子界，应在需要其他指数时另行核对。

\ProseParagraphBreak

全体输入的范数收敛反过来也会强迫一致有界性。假设每个 \(B_R^\delta\) 已在 \(L^p\)、\(1\leq p<\infty\) 上连续定义，并且对每个 \(f\) 都有 \(B_R^\delta f\to f\)。收敛控制大半径；在紧半径区间，缩放公式与 \(L^p\) 拉回的强连续性使每个输入的像连续，故也有界。第一卷定理 8.3.3 的一致有界原理于是给出 \(R\geq1\) 的统一算子界；半径缩放又说明各半径同指数范数相同，因此得到整个族的一致界。

\ProseParagraphBreak

这个推论需要“每个输入”与“已连续定义”两项条件。仅在 Schwartz 类上收敛没有给整个 Banach 空间的逐点有界性；一个不有界的形式乘子在测试函数上也可能表现出良好近似。因而证明一般输入的收敛时，先建立统一算子界，再用稠密性；反向用一致有界原理时，则必须已经拥有全部输入上的收敛。这两条路线互为检验，而不是可以循环引用的证明。

\ProseParagraphBreak

对于高维硬球，Fefferman 定理排除了非二指数的连续同指数算子，因此也排除了这一意义下对全部 \(L^p\) 输入的一致范数求和方案。它并不声称每个具体函数都失败，更不声称每个平滑测试函数失败。算子结论以全部输入和统一常数为量词，单个输入的偶然收敛不能否定它。

\subsection{核表示}
\label{b2:subsec:240103}

\begin{proposition}[Bochner--Riesz 核]\label{b2:prop:24-kernel}
对 \(\delta>-1\)，符号在球内可积，逆变换为连续有界径向函数，且
\[
 K_\delta(x)=\Gamma(\delta+1)\pi^{-\delta}
     |x|^{-d/2-\delta}J_{d/2+\delta}(2\pi|x|),\qquad x\ne0.
\]
在原点，
\[
 K_\delta(0)=\frac{\pi^{d/2}\Gamma(\delta+1)}
                    {\Gamma(\delta+d/2+1)}.
\]
\end{proposition}
\begin{proof}
由\cref{b2:thm:A2-bochner-riesz}在本卷 Fourier 归一化下的径向核公式直接得到；该处已用 Bessel 级数和 beta 积分给出完整推导。原点值取该级数首项，亦可直接积分球内符号。
\end{proof}
本章只把 \(\delta\ge0\) 当作通常平均；\(-1<\delta<0\) 的核恒等式用于说明解析族，并不意味着符号是 \(L^2\) 有界乘子。

% !TeX root = ../../Book2.tex
\section{核的渐近}
\label{b2:sec:2402}

核公式给出了一个特殊函数表达式，要判断可积性，还必须知道无穷远处的主项及误差。上界能够证明收敛，排除端点则需要下界。本节把球内积分化为两个一维端点的贡献，读出边界消失阶怎样增加空间衰减，再由同一渐近分别处理核范数、逼近恒等和极大控制。

\subsection{Bessel 表示}
\label{b2:subsec:240201}

除 Bessel 级数外，还有一个便于看清几何的切片表示。
设 \(d\geq2,\delta>-1\)，把空间点旋转为 \(x=(0,\ldots,0,r)\)。
频率写为 \((\eta,t)\in\mathbb R^{d-1}\times\mathbb R\)，则
\[
 \begin{aligned}
 K_\delta(r)
 &=\int_{-1}^1e^{2\pi\mathrm i rt}
       \int_{|\eta|^2<1-t^2}
          (1-t^2-|\eta|^2)^\delta\,d\eta\,dt.
 \end{aligned}
\]
所有换序都由球内幂权的绝对可积性保证。
内层令 \(\eta=(1-t^2)^{1/2}u\)，得到
\[
 \int_{|\eta|^2<1-t^2}(1-t^2-|\eta|^2)^\delta\,d\eta
 =c_{d,\delta}(1-t^2)^{\delta+(d-1)/2},
\]
其中用\cref{b2:prop:24-kernel}的原点积分在维数 \(d-1\) 中计算，
\[
 c_{d,\delta}
 =\frac{\pi^{(d-1)/2}\Gamma(\delta+1)}
        {\Gamma(\delta+(d+1)/2)}.
\]
故令 \(\alpha=\delta+(d-1)/2\)，有
\[
 K_\delta(r)=c_{d,\delta}
       \int_{-1}^1e^{2\pi\mathrm i rt}(1-t^2)^\alpha\,dt.
\]
当 \(d=1\) 时，直接从定义得到同一公式，此时 \(c_{1,\delta}=1\)。

\ProseParagraphBreak
横截面的尺度贡献 \((1-t^2)^{(d-1)/2}\)，原来的边界权再贡献
\((1-t^2)^\delta\)。两者相加，给出端点 \(t=\pm1\) 的消失阶
\(\alpha\)。区间内部的振幅光滑，且相位导数不为零，局部化后可
反复分部积分；真正决定主项的是两个端点。

\ProseParagraphBreak
低维模型可以同时检查常数与符号。先取 \(d=1\)，置 \(u=2\pi x\)。
零阶时直接积分得到
\[
 K_0(x)=\int_{-1}^1e^{\mathrm i u\xi}\,d\xi
       =\frac{2\sin u}{u},\quad K_0(0)=2.
\]
对一阶权 \(1-\xi^2\)，在有限区间上对 \(u\) 求导合法，因而
\[
 K_1(x)=\left(1+\frac{d^2}{du^2}\right)
              \frac{2\sin u}{u}
       =\frac{4(\sin u-u\cos u)}{u^3},\quad K_1(0)=\frac43.
\]
后一个原点值可由 Taylor 展开取得，也等于
\(\int_{-1}^1(1-\xi^2)\,d\xi\)。
当 \(u=2\pi\) 时，\(K_1(x)<0\)，所以即使正阶平均的核已可积，
它也不是一个非负的概率核。

\ProseParagraphBreak
在三维，零阶球的横截面积为 \(\pi(1-t^2)\)，于是
\[
 K_0(x)=4\pi\,\frac{\sin u-u\cos u}{u^3},
 \quad u=2\pi|x|,\quad K_0(0)=\frac{4\pi}{3}.
\]
它与一维一阶核只差常数，但绝对积分的径向体积元已经由 \(dr\)
变为 \(4\pi r^2\,dr\)。同一个函数形式在不同底空间有不同的
可积性，判断时不能省去体积增长。

\subsection{端点振荡渐近}
\label{b2:subsec:240202}

所需的一维工具是\cref{b2:lem:A2-endpoint-asymptotic}：
若 \(\beta>-1\)，\(\chi\) 在 \([0,\infty)\) 光滑且紧支撑，则
\[
 \int_0^\infty e^{-2\pi\mathrm i rs}s^\beta\chi(s)\,ds
 =\Gamma(\beta+1)(2\pi r)^{-\beta-1}
       e^{-\mathrm i\pi(\beta+1)/2}\chi(0)
       +O(r^{-\beta-2}).
\]
这里 \(r\to+\infty\)，复相位对应负号 Fourier 指数。
该结论不要求 \(\beta\) 为整数，因而能处理任意实消失阶。

\ProseParagraphBreak
回忆这个工具的证明机制，可以避免把端点公式误当成一次普通分部积分。
模型积分带上 \(e^{-s}\) 后为
\[
 \int_0^\infty e^{-2\pi\mathrm i rs}s^\beta e^{-s}\,ds
       =\Gamma(\beta+1)(1+2\pi\mathrm i r)^{-\beta-1}.
\]
将 \(\chi(s)\) 写成
\(\chi(0)e^{-s}+s\,c(s)\)，余项在零点多消失一阶。
在 \(0\leq s\lesssim r^{-1}\) 上直接积分，得到 \(O(r^{-\beta-2})\)；
在 \(s\asymp2^j/r\) 的剩余小环带上，取整数 \(N>\beta+2\)
次分部积分，得到
\[
 C r^{-\beta-2}2^{j(\beta+2-N)}.
\]
这是可求和的几何级数，远离零点的部分则有任意阶衰减。
这个分解既利用了端点的尺度，又没有对奇异振幅在零点作非法求导。
附录中的证明同时控制了截断导数与无穷远的指数尾部。

\ProseParagraphBreak
现在固定 \(d\geq1,\delta\geq0\)，将切片积分在两个端点与区间内部
作光滑分割。在 \(t=1\) 附近令 \(s=1-t\)，便有
\[
 (1-t^2)^\alpha=s^\alpha(2-s)^\alpha,\quad
 e^{2\pi\mathrm i rt}=e^{2\pi\mathrm i r}e^{-2\pi\mathrm i rs}.
\]
端点振幅在 \(s=0\) 的值为 \(2^\alpha\)。上述公式给出的正端点主项为
\[
 c_{d,\delta}\,2^\alpha\Gamma(\alpha+1)
       (2\pi r)^{-\alpha-1}
       e^{\,2\pi\mathrm i r-\mathrm i\pi(\alpha+1)/2}.
\]
负端点贡献是它的共轭，内部项经分部积分可忽略到所需阶数。
代入 \(\alpha=\delta+(d-1)/2\)，整理常数得到
\[
 \begin{aligned}
 K_\delta(r)
 &=a_\delta r^{-\delta-(d+1)/2}
       \cos\!\left(2\pi r-\frac{\pi(d+1+2\delta)}4\right)
       +O(r^{-\delta-(d+3)/2}),\\
 a_\delta&=\frac{\Gamma(\delta+1)}{\pi^{\delta+1}}>0.
 \end{aligned}
\]
这也与\cref{b2:prop:24-kernel}中 Bessel 阶数
\(\nu=d/2+\delta\) 的大半径渐近一致。

\ProseParagraphBreak
每增加一阶 \(\delta\)，空间主项便多衰减一个幂次。
在固定维数下，若 \(\delta\) 限制在 \([0,\infty)\) 的一个固定紧区间，
这里的余项常数还可以一致选取：端点分割固定，光滑因子
\((2-s)^\alpha\) 的有限阶导数一致有界；环带估计中的 \(N\) 取为
大于全部 \(\alpha+2\) 的同一个整数，幂函数积分的分母也远离零。
Gamma 系数在该参数紧集连续有界。这种局部参数一致性可用于研究
可积性阈值附近的范数变化，却不是允许阶数无界增长的渐近结论。

\ProseParagraphBreak
余弦零点并不会使主项在大部分半径上消失。令
\[
 \beta_\delta=\delta+\frac{d+1}{2},\quad
 \phi_\delta=\frac{\pi(d+1+2\delta)}4.
\]
在每个充分大的单位区间 \([n,n+1]\) 内，集合
\[
 A_n=\{r\in[n,n+1]:|\cos(2\pi r-\phi_\delta)|\geq3/4\}
\]
都有同一个正长度。余项比主振幅小一个 \(r^{-1}\) 因子，
因此当 \(n\) 足够大时，
\[
 |K_\delta(r)|\geq c_{d,\delta}r^{-\beta_\delta}
       \quad(r\in A_n).
\]
上界处理整个空间，下界只需占固定比例的这些径向区域。
下面将这两种用途分开应用。

\subsection{空间衰减}
\label{b2:subsec:240203}

\Needspace{6\baselineskip}
\begin{proposition}[核的有限指数可积性]\label{b2:prop:24-kernel-tail}
设 \(d\geq1,\delta\geq0\)，\(K_\delta\) 为
\cref{b2:def:24-bochner-riesz}的单位半径核。
对每个 \(1\leq q<\infty\)，有
\[
 K_\delta\in L^q(\mathbb R^d)
 \quad\Longleftrightarrow\quad
 q\left(\delta+\frac{d+1}{2}\right)>d.
\]
核在全空间有界；上述有限指数条件在等号时不成立。
\end{proposition}
\begin{proof}
核在有界集合上连续有界。远场上界给出
\[
 \int_{|x|>1}|K_\delta(x)|^q\,dx
 \leq C\int_1^\infty r^{d-1-q\beta_\delta}\,dr,
 \quad \beta_\delta=\delta+(d+1)/2.
\]
当 \(q\beta_\delta>d\) 时收敛。
反过来，在前述 \(A_n\) 上使用下界及极坐标，
\[
 \int_{|x|>1}|K_\delta(x)|^q\,dx
       \geq c\sum_{n\geq n_0}n^{d-1-q\beta_\delta}.
\]
这里 \(r\asymp n\)，且每个 \(A_n\) 的长度都有固定正下界。
若 \(q\beta_\delta\leq d\)，这个级数发散；临界等号时就是调和级数。
\end{proof}

这只判断一个核函数的 \(L^q\) 性质。乘子在 \(L^p\) 上有界是对
全部输入的要求，两者不能互换。核的 \(L^1\) 可积性能够通过 Young
提供所有指数的充分条件，但不是每个中间指数的必要条件。

\Needspace{6\baselineskip}
\begin{theorem}[可积核范围]\label{b2:thm:24-integrable}
设 \(d\geq1,\delta\geq0\)，\(B_R^\delta,K_\delta\)
如\cref{b2:def:24-bochner-riesz}。则
\[
 K_\delta\in L^1(\mathbb R^d)
       \quad\Longleftrightarrow\quad \delta>\frac{d-1}{2}.
\]
在这一范围，算子由 \(K_{\delta,R}*f\) 定义于所有
\(L^p(\mathbb R^d)\)、\(1\leq p\leq\infty\)，并有
\[
 \sup_{R>0}\|B_R^\delta f\|_p
       \leq\|K_\delta\|_1\cdot\|f\|_p.
\]
对 \(p<\infty\)，\(B_R^\delta f\to f\) 于 \(L^p\)；
对有界一致连续函数，收敛在全空间一致范数中成立。
\end{theorem}
\begin{proof}
在上一命题中取 \(q=1\)，得到准确的核可积阈值。
缩放不改变核的 \(L^1\) 范数，Young 不等式遂给出算子界。
这个卷积与 Schwartz 类上的 Fourier 定义一致，有限指数上的延拓
因而相容；无穷指数也与 \(L^1\) 伴随约定一致，因为核实值且为偶函数。

还须核对逼近恒等的总质量。符号与核均可积，Fourier 反演表明
\(\widehat K_\delta=m_{\delta,1}\)。二者均在原点连续，故
\[
 \int_{\mathbb R^d}K_\delta(y)\,dy=m_{\delta,1}(0)=1.
\]
对任何 \(\rho>0\)，缩放后的绝对尾质量满足
\[
 \int_{|x|>\rho}|K_{\delta,R}(x)|\,dx
       =\int_{|y|>R\rho}|K_\delta(y)|\,dy\longrightarrow0.
\]
这些是可以变号的逼近恒等条件。

更直接地，在 \(L^p\) 中有
\[
 B_R^\delta f-f
       =\int K_\delta(y)\,[f(\,\cdot-y/R)-f]\,dy.
\]
有限指数的平移连续性使每个固定 \(y\) 对应的差范数趋零，
且它至多为 \(2\|f\|_p\)。用 \(|K_\delta|\) 作可积控制，
Minkowski 与支配收敛即得范数极限。对有界一致连续输入，
同样使用平移的一致范数连续性，得到一致极限。
\end{proof}

在 \(L^\infty\) 上只陈述有界一致连续输入的收敛，不能把它扩张为
所有有界函数的一致逼近。核卷积可以平滑一个跳跃，却不能在一致范数
中把跳跃的两侧同时逼近到任意精度。后面的习题将把这个边界层算出来。
这里也没有使用核的非负性；总积分为一足以进入平移差公式。

\ProseParagraphBreak
可积阈值以上的径向衰减还控制全部半径的极大族。
设 \(\varepsilon=\delta-(d-1)/2>0\)，则
\[
 |K_\delta(x)|\leq C(1+|x|)^{-d-\varepsilon}.
\]
因而本问题满足\cref{b2:prop:17-approximate-identity}的具体核假设。
为看清严格正余量的用途，可以直接将卷积绝对值拆成小球与二进环带：
对 \(R>0\)，小球 \(|y|\leq R^{-1}\) 的贡献至多为 \(CMf(x)\)；
在
\(2^{j-1}R^{-1}<|y|\leq2^jR^{-1}\)、\(j\geq1\)，核大小与球体积相乘后
给出 \(C2^{-j\varepsilon}\)。于是
\[
 |B_R^\delta f(x)|
 \leq C\left(1+\sum_{j\geq1}2^{-j\varepsilon}\right)Mf(x).
\]

\Needspace{6\baselineskip}
\begin{corollary}[极大控制与 Lebesgue 点求和]\label{b2:cor:24-maximal}
设 \(d\geq1\)、\(\delta>(d-1)/2\)、\(1\leq p<\infty\)、
\(f\in L^p(\mathbb R^d)\)。采用核积分的代表，令
\[
 B_*^\delta f(x)\coloneqq\sup_{R>0}|B_R^\delta f(x)|.
\]
则 \(B_*^\delta f\leq C_{d,\delta}Mf\) 几乎处处。
它在 \(p>1\) 时强型有界，在 \(p=1\) 时为弱 \((1,1)\) 型。
在 \(f\) 的每个 Lebesgue 点 \(x\)，有
\[
 B_R^\delta f(x)\longrightarrow f(x)\quad(R\to\infty).
\]
\end{corollary}
\begin{proof}
前面的环带估计在 \(Mf(x)<\infty\) 处同时保证绝对收敛和极大上界。
对固定这样的 \(x\)，半径在任意紧区间内变化时，核函数连续，
并有与区间有关的共同可积控制；可用相同的环带和来作支配收敛。
所以 \(R\mapsto B_R^\delta f(x)\) 连续，上确界可取遍正有理半径，
从而是可测函数。极大函数的强、弱型界给出相应算子界。

若 \(x\) 为 Lebesgue 点，则 \(Mf(x)<\infty\)：小球平均由
Lebesgue 点性质控制，大球平均由 \(L^p\) 范数和 Hölder 控制，
其余半径留在紧区间。将卷积减去 \(f(x)\)，在每个环带上改用
\(|f(x-y)-f(x)|\) 的球平均。固定 \(j\) 时，球半径 \(2^j/R\to0\)，
该平均趋零；所有这些平均又被 \(Mf(x)+|f(x)|\) 控制。
由于 \(\sum_j2^{-j\varepsilon}<\infty\)，先控制级数尾，再令 \(R\to\infty\)，
就得到点态极限。
\end{proof}

如果把 \(\delta\) 降到核可积阈值，环带权不再可和，上述证明停止。
某个固定半径算子仍可能在中间指数上有界，但那并不自动控制极大族。
范数收敛、某个子列的几乎处处收敛与所有半径的几乎处处收敛，
需要各自的依据。

\ProseParagraphBreak
最后把两个端点的常数算清。对任意 \(K\in L^1(\mathbb R^d)\)，
卷积算子在 \(L^1\) 与 \(L^\infty\) 上的准确范数都是 \(\|K\|_1\)。
Young 已给出上界。在 \(L^1\) 上，取非负积分一的光滑近似恒等族
\(\rho_n\)，则
\[
 \|\rho_n\|_1=1,\quad K*\rho_n\to K\text{ 于 }L^1,
\]
这给出反向范数不等式。在 \(L^\infty\) 上，取
\[
 f(-y)=
 \begin{cases}
 \overline{K(y)}/|K(y)|,&K(y)\neq0,\\
 0,&K(y)=0.
 \end{cases}
\]
于是 \(\|f\|_\infty\leq1\) 且 \((K*f)(0)=\|K\|_1\)。
因为 \(K\) 的 \(L^1\) 平移连续，
\[
 |(K*f)(x+h)-(K*f)(x)|
       \leq\|K(\,\cdot+h)-K\|_1\cdot\|f\|_\infty,
\]
输出是连续函数，零点的值也给出本性上确界的下界。

\ProseParagraphBreak
应用于本章的可积核，得到
\[
 \|B_R^\delta\|_{1\to1}
       =\|B_R^\delta\|_{\infty\to\infty}
       =\|K_\delta\|_1>1.
\]
最后一个严格不等式来自振荡渐近：在无穷多径向区间内，
余弦主项为负且支配余项，所以核在一个正测度集合上为负。
核实值且总积分为一，因此
\(\int|K_\delta|=1+2\int(K_\delta)_- >1\)。
热核和 Poisson 核的非负性会给端点压缩；Bochner–Riesz 的总质量
归一化不具有同样的结论。

\ProseParagraphBreak
一维的普通指数范围至此可以完整辨认。零阶区间投影在
\(1<p<\infty\) 上由 Hilbert 变换有界，正阶再由
\cref{b2:prop:24-order-monotonicity}取得同一范围；
\(p=1,\infty\) 则要求正阶，核可积性给出充分性，下一节的尾部检验
给出必要性。高维中，仅靠这些事实仍留下大量中间指数，
这就是 Bochner–Riesz 问题的内容。

% !TeX root = ../../Book2.tex
\section{Bochner–Riesz 问题}
\label{b2:sec:2403}

可积核条件给出所有指数的范围，却未必是固定 \(p\) 的最小平滑量。中心指数 \(p=2\) 无需边界光滑；离开中心后，尾部与多方向集中逐渐成为障碍。

\subsection{\texorpdfstring{\(L^2\)}{L2} 理论}
\label{b2:subsec:240301}

对每个 \(\delta\ge0\)，符号的本质上确界为 \(1\)，所以
\[
 \|B_R^\delta\|_{2\to2}=1,\qquad B_R^\delta f\to f\text{ in }L^2.
\]
收敛由 Plancherel 与频率侧支配收敛直接得到。
但 \(L^2\) 范数界不控制频率边界上的点态取值，也不自动控制 \(L^p\)、\(p\ne2\)。与限制问题的联系需要额外结构。

\ProseParagraphBreak
对于一维 \(\delta=0\)，球就是区间，Hilbert 变换给出所有 \(1<p<\infty\) 上的界。高维球不再是有限个半轴投影的乘积。矩形投影可以逐坐标应用 Hilbert 变换并用 Fubini 得到有界性；球边界的法向连续变化，出现了无限方向的组合。

% 扩写定位：b2:subsec:240301
平方理论还允许量化平均逼近的速度，而且显示有无边界消失阶会改变速度结构。
\begin{proposition}\label{b2:prop:24-l2-rate}
设 \(\delta>0\)、\(0<s\leq2\)，且 \(f\in L^2\)、\((-\Delta)^{s/2}f\in L^2\)。则
\[
 \|B_R^\delta f-f\|_2
 \leq C_\delta(2\pi R)^{-s}\|(-\Delta)^{s/2}f\|_2.
\]
若 \(\delta=0\)，则对任意 \(s>0\) 同一结论成立，常数可以取一。
\end{proposition}
\begin{proof}
对 \(a\geq0\)，初等均值估计给出
\[
 |1-(1-a)_+^\delta|\leq\max\{1,\delta\}\min\{1,a\}.
\]
对 \(0<\delta<1\)，可由 \(1-(1-a)^\delta\leq a\) 的等价凹性关系验证；对 \(\delta\geq1\)，在 \(0\leq a\leq1\) 用导数上界 \(\delta\)。令 \(a=|\xi|^2/R^2\)。当 \(0<s\leq2\)，\(\min\{1,a\}\leq a^{s/2}\)，Plancherel 给出第一项。零阶硬截断的误差只在 \(|\xi|>R\)，在那里 \(\mathbf1_{\{|\xi|>R\}}\leq(|\xi|/R)^s\)，故可用任意正阶。
\end{proof}
对于 \(\delta>0\)，在一个固定紧频带内有
\(R^2[1-(1-|\xi|^2/R^2)^\delta]\to\delta|\xi|^2\)。
若非零 \(\widehat f\) 支于该频带，支配收敛于是给出
\[
 R^2(B_R^\delta f-f)\longrightarrow
 -\frac{\delta}{4\pi^2}(-\Delta)f
 \quad\text{于 }L^2.
\]
因此通常不能仅靠输入具有更多光滑性把第一项误差提高到 \(R^{-s}\)、\(s>2\)。要提高阶数，应组合不同平均以消去 Taylor 首项，或者采用在低频附近恰好等于一的另一种截断。频率边界的平滑与低频的逼近阶数是两个不同因素。

\ProseParagraphBreak

这个例子也说明“越平滑，逼近必然越快”不是无条件原则。硬截断在固定紧频带上最终误差为零；正阶平均虽然在球边界更光滑，却在球内部偏离一，因而留下低阶 Taylor 误差。硬截断的困难主要是其他 \(L^p\) 的一致有界性和逐点求和；平方逼近本身则十分直接。

\subsection{临界指标}
\label{b2:subsec:240302}

\begin{proposition}[尾部必要条件]\label{b2:prop:24-necessary}
设 \(1\le p\le\infty\)、\(\delta\ge0\)。若 \(B_1^\delta\) 在 \(L^p\) 有界，则在 \(p\ne2\) 的相关端点上必须满足
\[
 \delta>d\left|\frac1p-\frac12\right|-\frac12.
\]
当右侧为负数时，该条件不提供比 \(\delta\ge0\) 更强的限制。
\end{proposition}
\begin{proof}
选 \(\widehat f\in C_c^\infty\)，在闭单位球附近等于 \(1\)。则 \(f\) 为 Schwartz 函数，且 \(B_1^\delta f=K_\delta\)。
核渐近中的固定比例下界说明 \(K_\delta\in L^p\)、\(p<\infty\)，必须有
\(p[\delta+(d+1)/2]>d\)，即
\(\delta>d/p-(d+1)/2\)。
符号为实值，算子自伴随；若在 \(1<p<\infty\) 上有界，便也在 \(p'\) 上有界，对 \(p'\) 使用同一条件，合并得到结论。
当 \(p=1\)，直接得到 \(\delta>(d-1)/2\)。当 \(p=\infty\)，若有界，则对 \(L^1\cap L^2\) 的配对与 \(L^\infty\) 估计给出 \(L^1\) 界，故同样得到端点条件。对于 \(p=2\)，该检验给出 \(\delta>-1/2\)，所有 \(\delta\ge0\) 都满足。
\end{proof}
严格性来自径向调和级数；没有对尾部作精确下界，就不能宣称阈值等号也失败。这个检验并未排除临界函数
\(\max\{d|1/p-1/2|-1/2,0\}\) 中为零的内部指数上的硬截断；下一小节的深层几何定理才排除高维硬球。

\subsection{Bochner–Riesz 猜想}
\label{b2:subsec:240303}

经典 Bochner–Riesz 目标是：对 \(d\ge2\)、\(1<p<\infty\)，在
\[
 \delta>\max\left\{d\left|\frac1p-\frac12\right|-\frac12,\ 0\right\}
\]
时取得 \(L^p\) 有界性。这里严格的正平滑量将硬球情形排除在目标外。维数、端点形式、径向输入限制以及局部版本各是不同问题。本章以经典目标组织证明，不把未经本项目核验的当代进展范围记作定理。

\ProseParagraphBreak

下面采用\term[fefferman-qiuchengzi-dingli]{Fefferman 球乘子定理}[Fefferman ball multiplier theorem]的下述版本。

\begin{theorem}[深层结构定理：Fefferman 球乘子定理]\label{b2:thm:24-fefferman}
在 \(d\ge2\) 时，单位球特征函数为 \(L^p(\mathbb R^d)\) 乘子、
\(1<p<\infty\)，当且仅当 \(p=2\)。
\end{theorem}
\begin{proofsketch}
假设球乘子在某个非二指数有界。对球作放大、频移并取极限，配合随机符号，可推出各方向半空间投影的一致向量值界。Besicovitch 型细矩形构造使输入的平方函数质量分散，而输出在一个小测度集合中累积；对 \(p>2\) 这与向量值界矛盾。对偶性处理 \(p<2\)，切片传递处理高维。
完整的几何矩形构造及传递论证在此略去。正式原文为
\cite[Theorem 1; Section 2, pp.~330--336]{Fefferman1971Ball}。
该原文的出版扫描已核验；本段是证明路线摘要，不作为一个完整反例构造使用。
\end{proofsketch}
这个定理并不说任何正的 \(\delta\) 都足够。球乘子的失败与正平滑后的范围之间，仍隔着一个依赖维数与 \(p\) 的估计问题。

\subsection{已知范围的结构性概览}
\label{b2:subsec:240304}

\begin{theorem}[可直接证明的插值范围]\label{b2:thm:24-interpolation-range}
对 \(1<p<\infty\)，若
\[
 \delta>(d-1)\left|\frac1p-\frac12\right|,
\]
则 \(B_R^\delta\) 在 \(L^p\) 上一致有界，且随 \(R\to\infty\) 在 \(L^p\) 中逼近恒等。
\end{theorem}
\begin{proof}
取 \(\sigma>(d-1)/2\)。对符号族 \((1-|\xi|^2)_+^z\)，在
\(\operatorname{Re}z=0\) 线上有 \(L^2\) 范数 \(1\)。
在 \(\operatorname{Re}z=\sigma\) 线上，复阶 Bessel 方程与级数给出
\[
 |K_{\sigma+it}(x)|\le C e^{C(1+t^2)}
                (1+|x|)^{-\sigma-(d+1)/2}.
\]
核在大半径的界使用\chapref{b2:ch:23}复阶 ODE 能量证明（阶数实部改变，但仍固定），小半径使用 Gamma 递推和级数。右侧可积，故有 \(L^1\) 与 \(L^\infty\) 界，增长至多 Gaussian 型。
给解析族乘 \(e^{\kappa z^2}\) 消去竖线增长，应用解析插值：
在 \(\delta=\theta\sigma\) 上，
\(|1/p-1/2|=\theta/2\)。
因此任意 \(\delta>(d-1)|1/p-1/2|\) 都可选一个这样的 \(\sigma\)。
缩放给出所有半径的一致界，上一节的稠密性论证给出逼近。
\end{proof}
这是一个明确已证范围，一般窄于经典目标。它的证明只结合 \(L^2\) 与可积核端点，没有利用球面限制的额外改进。比较两条临界直线，能看清深层理论应在哪个区间增加信息。

已经证明的一致乘子界也能与输入正则性配合，给出其他 \(L^p\) 中的逼近速度。
\begin{proposition}\label{b2:prop:24-lp-rate}
设 \(1<p<\infty\)、\(\delta>0\)，且 \(B_R^\delta\) 在该 \(L^p\) 上一致有界。则对 \(0<s\leq2\)、\(R\geq1\)，
\[
 \|B_R^\delta f-f\|_p\leq C R^{-s}\|f\|_{H^{s,p}}.
\]
\end{proposition}
\begin{proof}
取 \(\chi\in C_c^\infty\)，在 \(|\xi|\leq1/4\) 上为一，支撑于 \(|\xi|<1/2\)。把输入分为 \(\chi(\xi/R)\) 的低频部与其补的高频部。记 \(g=J^sf\)。高频部的乘子
\[
 R^s(1-\chi(\xi/R))(1+4\pi^2|\xi|^2)^{-s/2}
\]
满足统一 Mihlin 界，因为它支于 \(|\xi|\geq R/4\)，在该处 \(R\leq4|\xi|\)。因此高频输入的 \(L^p\) 范数不超过 \(CR^{-s}\|g\|_p\)。对它再用 \(I-B_R^\delta\) 的一致界即可。

低频误差的归一化乘子为
\[
 R^s[1-(1-|\xi|^2/R^2)^\delta]\,
 \chi(\xi/R)(1+4\pi^2|\xi|^2)^{-s/2}.
\]
在这片频率中边界离支集有固定距离；把方括号写成
\((|\xi|^2/R^2)b_\delta(|\xi|^2/R^2)\)，其中 \(b_\delta\) 及所需导数一致有界。对 \(|\xi|\leq1\)，因 \(s\leq2\)、\(R\geq1\)，前因子 \(R^{s-2}\) 有界；对 \(1\leq|\xi|\leq R/2\)，因子可写为 \((|\xi|/R)^{2-s}\) 乘一个零阶有界光滑量。逐次链式求导便得到统一的尺度导数界。Mihlin 给出低频误差上界，合并两部分证明命题。
\end{proof}
这个论证只要求本章某条途径已经给出一致乘子界，并不把充分条件误当成必要条件。它也把两种误差分开：高频尾部由输入的光滑性控制，低频误差由平均符号在原点的 Taylor 阶数控制。若进一步构造消掉低阶 Taylor 系数的组合平均，后一限制可以改变，但需要重新核对组合的统一算子界。

\ProseParagraphBreak

例如，对固定 \(d,p\) 选择处于已证插值范围的 \(\delta\)，一个 \(W^{1,p}\) 输入至少有 \(R^{-1}\) 的范数逼近速度，一个 \(W^{2,p}\) 输入有 \(R^{-2}\) 的速度。这里只是范数误差，不能把它直接变成全参数逐点误差率；后者需要位置上的极大控制或更强的正则性与嵌入。

\ProseParagraphBreak

必要直线与插值直线的距离可以用一个数值模型说明。在 \(d=3\)、\(p=4\) 中，尾部条件要求 \(\delta>1/4\)，本章直接证明的插值范围要求 \(\delta>1/2\)。例如 \(\delta=3/4\) 在已证范围内，\(\delta=1/8\) 被必要条件排除，而 \(\delta=3/8\) 不能仅靠这两项结果判断。说它“没有被反例排除”与说它“已证明有界”有明确区别。

\ProseParagraphBreak

指数对偶把 \(p=4\) 与 \(p=4/3\) 的范围对应起来，因为球符号为实且算子在测试类上自伴随。缩放又使每个已证同指数界自动对所有半径具有同一常数。因此这里留下的空缺是平滑量与输入指数的估计空缺，而不是半径依赖尚未处理。若一篇文献只给固定空间球内的局部输出或径向输入，其命题还不能直接填入这个全空间、一般输入的空缺。

\ProseParagraphBreak

在 \(p\) 接近二时，尾部必要条件的右侧可能为负，此时它只要求原来已经约定的 \(\delta\geq0\)，看起来没有门槛。然而高维 \(\delta=0\) 已由球乘子定理排除所有非二指数。这个额外的几何障碍来自多方向叠加，并不是单个核尾部的一条更仔细积分能看出来的。正阶问题是否能在接近二的范围内降到任意小的平滑量，则要依赖更深的曲率和限制估计。

\ProseParagraphBreak

对角线上的 \(p=2\) 则没有这些困难：符号有界已经给出范数一，且所有 \(\delta\geq0\) 都随固定阶数的半径增加强收敛到恒等。这个特殊点是各种插值证明的可靠起点，但它不控制其他指数上的几何组合。把一个证明完全归结为“平方理论加插值”时，应能够说清另一个端点来自哪里，以及为什么它有足够强的范围和统一常数。

% !TeX root = ../../Book2.tex
\section{与限制问题的关系}
\label{b2:sec:2404}

球面限制控制薄壳上的 Fourier 质量，Bochner–Riesz 符号则把不同薄壳加权求和。两者共享曲率与频率集中，但混合指数界和同指数乘子界不同，不能仅凭一个球面定理直接宣称整个乘子猜想成立。

\subsection{球面 Fourier 变换}
\label{b2:subsec:240401}

极坐标给出层分解
\[
 B_1^\delta f(x)=\int_0^1(1-r^2)^\delta r^{d-1}
        \int_{S^{d-1}}e^{2\pi\mathrm i r x\cdot\omega}
                  \widehat f(r\omega)\,\mathrm d\sigma(\omega)\,\mathrm dr.
\]
在 Schwartz 输入上换序合法，因为球内符号可积。内层正是半径 \(r\) 的球面限制再延拓。
球面本身的 Fourier 变换在径向尾部衰减，径向权的边界消失再增加 \(\delta+1\) 的振荡增益；这与核渐近的总阶数一致。

\ProseParagraphBreak
若只把内层绝对值估计后对 \(r\) 积分，就会丢失不同半径之间的抵消。某些混合指数估计仍可由这一办法直接证明，但锐利同指数乘子问题往往需要保留空间与径向的振荡。

\subsection{限制估计}
\label{b2:subsec:240402}

令 \(p_*=2(d+1)/(d+3)\)、\(q_*=p_*'\)。对薄壳
\(\Omega_\epsilon=\{1-\epsilon<|\xi|<1+\epsilon\}\)，球面限制的缩放给出
\[
 \|\mathbf1_{\Omega_\epsilon}\widehat f\|_2
       \le C\epsilon^{1/2}\|f\|_{p_*},\qquad0<\epsilon<1/2.
\]
证明是将平方范数写成半径积分，逐半径应用限制界；半径在固定紧区间，其缩放常数一致。
与对偶估计复合，或直接逐半径限制再延拓，得到
\[
 \|\mathcal F^{-1}(\mathbf1_{\Omega_\epsilon}\widehat f)\|_{q_*}
       \le C\epsilon\|f\|_{p_*}.
\]
第一式控制一个频率 \(L^2\) 范数；第二式是 \(L^{p_*}\to L^{q_*}\) 的混合指数界，不是 \(L^{p_*}\to L^{p_*}\)。

\ProseParagraphBreak
对薄壳内一个非负径向权 \(w(r)\)，同样有算子范数至多
\(C\int|w(r)|\,\mathrm dr\)。在边界距离约 \(2^{-j}\) 的壳层上，Bochner–Riesz 权大小约 \(2^{-j\delta}\)，得到混合指数贡献 \(C2^{-j(\delta+1)}\)。因此在这条混合指数线上，壳层求和对 \(\delta>-1\) 可收敛；这与负参数不属于 \(L^2\) 有界乘子并不矛盾。

把半径趋零的缩放也保留下来，可以完整证明一个全平均的混合指数界。
\begin{proposition}\label{b2:prop:24-mixed-bound}
设 \(d\geq2\)、\(p_*=2(d+1)/(d+3)\)、\(q_*=p_*'\)。对任意 \(\delta>-1\)，以球内可积符号在 Schwartz 类上定义的 \(B_1^\delta\) 有唯一连续延拓
\[
 B_1^\delta:L^{p_*}(\mathbb R^d)\longrightarrow L^{q_*}(\mathbb R^d).
\]
\end{proposition}
\begin{proof}
记 \(T_r f(x)=\int_{S^{d-1}}e^{2\pi\mathrm irx\cdot\omega}
\widehat f(r\omega)\,\mathrm d\sigma(\omega)\)。设 \(f_r(y)=f(y/r)\)，则
\(\widehat f(r\omega)=r^{-d}\widehat{f_r}(\omega)\)，因而
\[
 T_r f(x)=r^{-d}(ERf_r)(rx),\qquad
 \|T_r f\|_{q_*}\leq C r^{d/p_*-d/q_*-d}\|f\|_{p_*}.
\]
层分解中的体积权 \(r^{d-1}\) 使最后指数成为
\(d/p_*-d/q_*-1=(d-1)/(d+1)>-1\)。
于是 Minkowski 与限制端点给出
\[
 \|B_1^\delta f\|_{q_*}\leq C\|f\|_{p_*}
 \int_0^1(1-r^2)^\delta r^{(d-1)/(d+1)}\,\mathrm dr.
\]
积分在 \(r=0\) 的幂次可积，在 \(r=1\) 可积恰需 \(\delta>-1\)。稠密性完成延拓。
\end{proof}
这条命题包括负参数，但改变了输出指数，因此不与负参数符号在边界无界矛盾。对任意有界同指数 \(L^2\) 乘子，Plancherel 需要符号本性有界；这里使用的是较强可积性输入换取较弱的全空间输出，根本没有提出平方到平方的界。

\ProseParagraphBreak

半径 \(R\) 的相同混合界还会带缩放因子
\(R^{d(1/p_*-1/q_*)}\)，不像同指数乘子范数那样独立于 \(R\)。如果在转述中删除这个因子，便会把不同类型的尺度关系混在一起。对固定单位半径的估计，缩放检查仍然有价值：它可以迅速发现体积权、曲面测度与 Fourier 归一化中的遗漏。

\subsection{乘子估计}
\label{b2:subsec:240403}

若要求同指数 \(L^p\)，上述输出 \(q_*>p_*\) 不能通过无限测度空间上的包含关系改成 \(p_*\)。一种方法是在空间中再按尺度局部化核，使用限制界控制一部分项、核尾部控制另一部分项，然后在壳层间建立可求和的界。这个过程涉及局部球的体积因子与几乎正交，不是形式上的指数替换。

\ProseParagraphBreak
\chapref{b2:ch:19}平方函数处理互相分离的频率环带，而本章的壳层都贴近同一个球面。它们在频率上可不交，却仍可能在空间中沿相近方向高度相干。普通 Plancherel 正交只给 \(p=2\)，其他指数上的组合需额外曲率或多线性控制。
现代解耦和限制方法正是组织这种相干的重要工具，具体结果必须按已出版定理核对。

\subsection{曲率}
\label{b2:subsec:240404}

球、椭球与平面多面体提供三种可比较几何。固定非退化线性换元把椭球变成球，因此同指数有界性范围及算子范数相同；各方向导数的估计则可能依赖换元矩阵。矩形可拆成各坐标区间投影，硬截断在 \(1<p<\infty\) 上有界。球面法向连续变化，硬截断则受 Fefferman 定理阻碍。

\ProseParagraphBreak
曲率既能改善单个振荡核衰减，也能产生大量不同方向的集中管。前一种效应使正平滑平均有较好的核估计，后一种效应使硬球乘子失败。它们并不矛盾，因为核尾部的单输入检验和多方向向量值检验衡量不同机制。
阅读乘子范围时，应明确平滑量、输入是否径向、空间是否局部、半径是否统一及是否有尺度损失；这些条件决定一个结论能否被用于普通 \(L^p(\mathbb R^d)\) 平均。

曲率条件的比较还可以作成实际的变换。若 \(A\) 为可逆矩阵，椭球符号 \((1-|A\xi|^2)_+^\delta\) 的算子由 Fourier 换元与空间线性拉回共轭于球平均。拉回在 \(L^p\) 的范数因子为 \(|\det A|^{1/p}\) 或其逆，共轭时两者相消；同指数算子范数因此与标准球相同。若改变的同时还使用不同的空间欧氏导数范数，其正则性常数才会依赖矩阵的条件数。乘子有界性与导数估计的矩阵依赖不能混为一谈。

\ProseParagraphBreak

多面体的情况并非只由“没有曲率”决定。一个固定有界凸多面体是有限个半空间的交，每个半空间乘子可通过旋转、调制和一维 Hilbert 投影在 \(1<p<\infty\) 上实现；这些 Fourier 乘法算子彼此交换，有限次复合就得到多面体指示函数乘子。其界可能随面数积累，所以用越来越多面的多面体逼近球，并没有产生球乘子的一致界。几何对象的逐点逼近若缺少一致算子常数，不能用于转移有界性。

\ProseParagraphBreak

球面边界有连续的一族法向方向，Fefferman 的反例正是把越来越多方向一起测试。另一方面，对一个固定方向的单个驻相片，非零曲率又能提供核衰减。分析某个新频率区域时，应同时询问局部片的衰减怎样、总共有多少片、不同片怎样组合；只有第一项回答得很好，仍可能在全局组合时失败。限制理论、多线性横截性和解耦就是为这种组合问题提供额外结构。

\ProseParagraphBreak

本章保留的必要条件与充分范围因此承担不同职责。尾部检验排除不可能的指数，可积核与解析插值提供完整可复核的充分区间，限制端点给出另一类混合指数界，深层球乘子定理解释零阶硬截断的特殊障碍。对某个尚未覆盖的正阶参数，这些结果共同界定问题所在，但没有自动给出结论。阅读研究文献中的改进时，应把其参数和模型放回这些已经明确的坐标中，才能判断它究竟改善了哪一段范围。


\begin{chaptersummary}
核可积的阈值为 \(\delta>(d-1)/2\)，中心指数 \(p=2\) 则允许硬截断。固定 \(p\) 的尾部检验与解析插值分别给出必要和充分直线，二者之间正是曲率与薄壳理论要处理的区间。限制估计首先产生混合指数界；同指数平均还须额外组合机制。
\end{chaptersummary}

\BookExercises

% 自拟：围绕本章已经建立的工具，改变条件或组合方法。
\begin{exercise}\label{b2:ex:24-new-1}
在 \(d=1,\delta=0\) 时核公式应化成什么？检查所有常数。
\end{exercise}
\begin{solution}
\(J_{1/2}(2\pi r)=\sin(2\pi r)/(\pi\sqrt r)\)，所以 \(K_0(x)=\sin(2\pi x)/(\pi x)\)，原点极限为2。这正是 \(\mathbf1_{[-1,1]}\) 的逆 Fourier 变换。
\end{solution}

% 自拟：围绕本章已经建立的工具，改变条件或组合方法。
\begin{exercise}\label{b2:ex:24-new-2}
证明 \(K_\delta(0)=\pi^{d/2}\Gamma(\delta+1)/\Gamma(\delta+d/2+1)\)。
\end{exercise}
\begin{solution}
原点处逆变换等于球内符号积分。用球面面积 \(2\pi^{d/2}/\Gamma(d/2)\) 乘径向积分，再令 \(t=r^2\)，径向积分为 \(\tfrac12 B(d/2,\delta+1)\)。Gamma 因子消去后得到公式。
\end{solution}

% 自拟：围绕本章已经建立的工具，改变条件或组合方法。
\begin{exercise}\label{b2:ex:24-new-3}
对 \(d=3\) 写出核可积阈值、尾部必要指标和本章插值充分指标。
\end{exercise}
\begin{solution}
核可积要求 \(\delta>1\)。必要指标为 \(\delta>3|1/p-1/2|-1/2\)（若右侧为负则无额外限制），充分指标为 \(\delta>2|1/p-1/2|\)。在 \(p=2\) 上，另外有直接的 \(\delta\ge0\) 有界性。三个陈述的用途不同。
\end{solution}

% 自拟：围绕本章已经建立的工具，改变条件或组合方法。
\begin{exercise}\label{b2:ex:24-new-4}
解释核渐近中的余弦零点为什么不能挽救 \(L^1\) 临界指标。
\end{exercise}
\begin{solution}
每个足够大的固定长度径向区间内，余弦绝对值大于 \(1/2\) 的子集有统一正长度。余项比主项低一个幂次，故在这些子集上仍有主项的一半下界。临界径向绝对积分的各区间贡献为 \(c/n\)，因此发散。
\end{solution}

% 自拟：围绕本章已经建立的工具，改变条件或组合方法。
\begin{exercise}\label{b2:ex:24-new-5}
已知单位半径平均在 \(L^p\)、\(1<p<\infty\) 上有界，证明所有半径的范数相同并得到 \(L^p\) 逼近。
\end{exercise}
\begin{solution}
缩放共轭公式中的两缩放范数相消，先得一个方向，再用逆缩放得相等。对频率紧支撑的 Schwartz 输入，随半径增大符号在支撑上趋1且导数趋零，得到 Schwartz 收敛。稠密性和统一范数界推广到所有 \(L^p\)。
\end{solution}

% 自拟：围绕本章已经建立的工具，改变条件或组合方法。
\begin{exercise}\label{b2:ex:24-new-6}
从球面 Stein–Tomas 定理推出薄壳 Fourier 质量界中的 \(\epsilon^{1/2}\) 因子。
\end{exercise}
\begin{solution}
极坐标把平方范数写成 \(\int_{1-\epsilon}^{1+\epsilon}r^{d-1}\|\widehat f(r\,\cdot)\|_{L^2(S)}^2\,\mathrm dr\)。逐半径限制界的缩放常数在该区间一致，积分长度为 \(2\epsilon\)，所以平方范数至多 \(C\epsilon\|f\|_{p_*}^2\)，取平方根即得。
\end{solution}

% 自拟：围绕本章已经建立的工具，改变条件或组合方法。
\begin{exercise}\label{b2:ex:24-new-7}
证明固定椭球频率截断的同指数乘子性质与球相同。
\end{exercise}
\begin{solution}
若椭球为 \(\{|A\xi|\le1\}\)、\(A\) 可逆，空间线性换元的 Fourier 公式将其算子与球截断共轭。两个换元算子的 \(L^p\) 范数分别含互逆 determinant 因子，因此有限性等价。于是 \(d\ge2\) 的硬椭球在非二指数仍不有界。
\end{solution}

% 自拟：围绕本章已经建立的工具，改变条件或组合方法。
\begin{exercise}\label{b2:ex:24-new-8}
为什么 \(L^{p_*}\to L^{p_*'}\) 的薄壳估计不能直接推出 \(L^{p_*}\) 球乘子界？给出函数空间上的具体障碍。
\end{exercise}
\begin{solution}
在 \(\mathbb R^d\) 上 \(L^{p_*'}\) 不包含 \(L^{p_*}\)。例如选择 \(a\) 满足 \(d/p_*'<a\le d/p_*\)，函数 \((1+|x|)^{-a}\) 属于前者而不属于后者。因而输出指数无法靠嵌入改变；还需空间局部化与额外估计。Fefferman 定理也直接排除硬球的非二同指数结论。
\end{solution}
